% Section 11.1: Tangent vectors as derivations; coordinate bases | Part: 2 | Depends on: sec:manifolds, sec:vectors, sec:ext | Status: todo | Original lines: 1863-1883
%!TEX root = ../../../main.tex
\subsection{Tangent vectors as derivations; coordinate bases}\label{sec:tangentvectors}
\begin{intuition}[{velocities}]
Let $p\in M$. If you move through $p$ along a curve $\gamma(t)$, your velocity $\dot\gamma(0)$ is an ``infinitesimal arrow'' at $p$. The set of all these velocities at $p$ is the \textbf{tangent space} $T_pM$, a flat $m$-dimensional vector space attached to $p$ (the ``tangent plane'' for a surface). Different points have different (a priori unrelated!) tangent spaces; that is the root of the difficulty with differentiating vector fields on a manifold in Section~\ref{sec:conn}. Covectors at $p$ live in the \textbf{cotangent space} $T^*_pM=(T_pM)^*$: they are the ``measuring devices'' for velocities.

We cannot describe a velocity on an abstract manifold as an arrow in some ambient $\R^N$; but we \emph{can} describe it by what it does to functions: a velocity $v$ at $p$ gives, for each smooth function $f$, the rate of change $v(f)$ of $f$ along the motion.
\end{intuition}

\begin{formal}[{tangent space as derivations; coordinate bases}]
A \emph{derivation at $p$} is a linear map $D:C^\infty(M)\to\R$ with $D(fg)=D(f)g(p)+f(p)D(g)$. The set of derivations is a vector space of dimension $m$, denoted $T_pM=\Der_p(M)$; its dual is $T^*_pM$. For a velocity of a curve $D_\gamma f=\frac{\dd}{\dd\varepsilon}f(\gamma(\varepsilon))|_{\varepsilon=0}$, and in $\R^m$ every derivation is a directional derivative $D_{p,v}f=\frac{\dd}{\dd\varepsilon}f(p+\varepsilon v)|_0$.

Given a chart $(U,\varphi)$ with coordinates $x^i$, define
\[
\partial_i|_p f:=\frac{\partial f}{\partial x^i}(p)=\frac{\dd}{\dd\varepsilon}\Big(f\circ\varphi^{-1}(\varphi(p)+\varepsilon e_i)\Big)\Big|_{\varepsilon=0},\qquad
\dd x^i|_p:\ \text{ the dual basis.}
\]
Then $\{\partial_i|_p\}$ is the \emph{coordinate basis} of $T_pM$ and $\{\dd x^i|_p\}$ the dual basis of $T_p^*M$, \ie $\dd x^i(\partial_j)=\delta^i_j$. The \emph{differential} of $f\in C^\infty(M)$ is the covector $\dd f_p(v)=v(f)$, with $\dd f=\partial_if\,\dd x^i$. Disjoint unions $TM=\bigsqcup_pT_pM$ and $T^*M$ are the \emph{tangent} and \emph{cotangent bundles}, manifolds of dimension $2m$.
\end{formal}
